% =============================================================================
%  TEMPLATE ARXiv REFINADO v1.0
%  Uso: submissao direta ao arXiv (http://arxiv.org)
%  Compilar: pdflatex -> pdflatex
% =============================================================================
\documentclass[12pt,a4paper]{article}

% --- Pacotes essenciais para arXiv ---
\usepackage[utf8]{inputenc}
\usepackage[T1]{fontenc}
\usepackage{amsmath,amssymb,amsthm}
\usepackage{graphicx}
\usepackage[margin=2.5cm]{geometry}
\usepackage[brazil]{babel}
\usepackage{xcolor}
\usepackage{hyperref}
\hypersetup{
    colorlinks=true,
    linkcolor=blue!60!black,
    citecolor=green!40!black,
    urlcolor=blue!60!black,
    breaklinks=true
}
% \usepackage{natbib}  % arXiv: usar thebibliography padrão
\usepackage{booktabs}
\usepackage{caption}
\usepackage{subcaption}
\usepackage{enumitem}
\usepackage{algorithm}
\usepackage{algpseudocode}

% --- Ambiente teoremas ---
\theoremstyle{plain}
\newtheorem{theorem}{Teorema}[section]
\newtheorem{lemma}[theorem]{Lema}
\newtheorem{corollary}[theorem]{Corol\'ario}

\theoremstyle{definition}
\newtheorem{definition}[theorem]{Defini\c{c}\~ao}
\newtheorem{example}[theorem]{Exemplo}

\theoremstyle{remark}
\newtheorem{remark}[theorem]{Observa\c{c}\~ao}

% --- Comandos auxiliares ---
\newcommand{\R}{\mathbb{R}}
\newcommand{\N}{\mathbb{N}}
\newcommand{\E}{\mathbb{E}}
\newcommand{\Var}{\operatorname{Var}}

% =============================================================================
\begin{document}

\title{T\'itulo do Artigo --- Vers\~ao Refinada para arXiv}

\author{Marcos A. Santos\thanks{Departamento de Computa\c{c}\~ao, Universidade Federal
do Cear\'a. E-mail: marcos.santos@exemplo.com},
        Ana C. Oliveira\thanks{Instituto de Ci\^encias Matem\'aticas,
Universidade de S\~ao Paulo. E-mail: ana.oliveira@exemplo.com}}

\date{\today}

\maketitle

% =============================================================================
\begin{abstract}
Este documento serve como template refinado para submiss\~ao ao arXiv.
Apresenta a estrutura padr\~ao de um artigo cient\'ifico com se\c{c}\~oes
bem definidas, ambientes matem\'aticos, tabelas, figuras, algoritmos e
refer\^encias no formato num\'erico. O c\'odigo-fonte est\'a preparado para
compila\c{c}\~ao direta com \texttt{pdflatex} e atende aos requisitos do
reposit\'orio arXiv.
\end{abstract}

% =============================================================================
\section{Introdu\c{c}\~ao}
\label{sec:intro}

A motiva\c{c}\~ao deste trabalho \'e demonstrar um template completo e
refinado para submiss\~ao de artigos ao reposit\'orio arXiv
\cite{arXiv2024Guide}. O arXiv aceita submiss\~oes em \LaTeX
desde que o c\'odigo-fonte seja compat\'ivel com as classes padr\~ao
(\texttt{article}, \texttt{amsart}, \texttt{revtex4-2}, etc.).

Um aspecto fundamental \'e a reprodu\c{c}\~ao confi\'avel do documento:
todas as fontes, pacotes e imagens devem estar inclu\'idos na submiss\~ao.
Pacotes n\~ao padr\~ao devem ser evitados ou inclu\'idos junto ao c\'odigo
\cite{arXivSubmitFAQ}.

% =============================================================================
\section{Fundamenta\c{c}\~ao Te\'orica}
\label{sec:fund}

\subsection{Formula\c{c}\~ao Matem\'atica}

Seja $\mathcal{X} \subseteq \R^n$ um espa\c{c}o de busca e
$f: \mathcal{X} \to \R$ uma fun\c{c}\~ao objetivo. O problema de
otimiza\c{c}\~ao consiste em encontrar

\begin{equation}
    x^* = \arg\min_{x \in \mathcal{X}} f(x)
    \label{eq:otim}
\end{equation}

\noindent sob um conjunto de restri\c{c}\~oes $g_i(x) \leq 0$ para
$i = 1, \dots, m$.

\begin{theorem}[Converg\^encia]
\label{thm:conv}
Sob as hip\'oteses de convexidade de $f$ e diferenciabilidade
com constante de Lipschitz $L$, o m\'etodo do gradiente descendente
com taxa de aprendizado $\eta < 2/L$ converge para o \'otimo global:
\begin{equation}
    \|x_{k+1} - x^*\|^2 \leq \left(1 - \frac{2\eta}{L}\right)
    \|x_k - x^*\|^2
    \label{eq:conv_rate}
\end{equation}
\end{theorem}

\begin{proof}
A prova segue diretamente da expans\~ao de Taylor de $f$ em torno de
$x_k$ e da propriedade de Lipschitz do gradiente \cite{Boyd2004Convex}.
\end{proof}

\subsection{An\'alise Estat\'istica}

Para um conjunto de $n$ observa\c{c}\~oes independentes $\{y_i\}_{i=1}^n$,
o estimador de m\'axima verossimilhan\c{c}a \'e dado por

\begin{equation}
    \hat{\theta}_{\text{MV}} = \arg\max_{\theta \in \Theta}
    \sum_{i=1}^n \log p(y_i \mid \theta)
    \label{eq:mv}
\end{equation}

\begin{table}[htbp]
\centering
\caption{M\'etricas de desempenho para diferentes configura\c{c}\~oes}
\label{tab:metricas}
\begin{tabular}{lccc}
\toprule
Configura\c{c}\~ao & Precis\~ao (\%) & Recall (\%) & $F_1$ \\
\midrule
Baseline       & 82.3 $\pm$ 1.2 & 79.8 $\pm$ 1.5 & 0.810 \\
Proposta       & 94.7 $\pm$ 0.8 & 93.1 $\pm$ 0.9 & 0.939 \\
Estado da arte & 95.2 $\pm$ 0.6 & 94.5 $\pm$ 0.7 & 0.948 \\
\bottomrule
\end{tabular}
\end{table}

% =============================================================================
\section{Metodologia}
\label{sec:metod}

O pipeline proposto segue tr\^es etapas principais, conforme descrito
no Algoritmo~\ref{alg:pipeline}.

\begin{algorithm}[htbp]
\caption{Pipeline de processamento}
\label{alg:pipeline}
\begin{algorithmic}[1]
\Require Conjunto de dados $\mathcal{D}$, par\^ametros $\Theta$
\Ensure Sa\'ida processada $\mathcal{Y}$
\State $\mathcal{X} \gets \text{Pr\'e-processamento}(\mathcal{D})$
\State $\mathcal{F} \gets \text{Extra\c{c}\~ao de caracter\'isticas}(\mathcal{X})$
\For{cada amostra $x_i \in \mathcal{X}$}
    \State $s_i \gets \text{Modelo}(\mathcal{F}_i, \Theta)$
    \If{$s_i > \tau$}
        \State $\mathcal{Y} \gets \mathcal{Y} \cup \{x_i\}$
    \EndIf
\EndFor
\State \Return $\mathcal{Y}$
\end{algorithmic}
\end{algorithm}

% =============================================================================
\section{Resultados}
\label{sec:result}

Os experimentos foram conduzidos em um conjunto de 10~000 amostras
com valida\c{c}\~ao cruzada $k$-fold ($k=5$). A Figura~\ref{fig:curva}
ilustra a curva de aprendizado do modelo proposto.

\begin{figure}[htbp]
\centering
% \includegraphics[width=0.7\textwidth]{curva_aprendizado.pdf}
\vspace{3cm}
\begin{minipage}{0.7\textwidth}
\centering
\fbox{\parbox{\textwidth}{\centering
    [Gr\'afico: Curva de perda vs. \'epocas]}}
\end{minipage}
\caption{Curva de aprendizado do modelo proposto.}
\label{fig:curva}
\end{figure}

% =============================================================================
\section{Discuss\~ao}
\label{sec:disc}

Os resultados indicam uma melhoria significativa em rela\c{c}\~ao
ao baseline (Tabela~\ref{tab:metricas}), com ganho de 12.4 p.p.
na m\'etrica $F_1$. A abordagem proposta supera o estado da arte
em 2 dos 3 datasets avaliados.

Limita\c{c}\~oes incluem a depend\^encia do hiperpar\^ametro $\tau$
e o custo computacional da etapa de pr\'e-processamento.

% =============================================================================
\section{Conclus\~ao}
\label{sec:conc}

Apresentamos um template refinado para submiss\~ao ao arXiv, com
estrutura completa de artigo cient\'ifico. O c\'odigo est\'a
preparado para compila\c{c}\~ao direta e pode ser adaptado para
diferentes \'areas do conhecimento.

Trabalhos futuros incluem a extens\~ao para m\'ultiplos datasets
e a otimiza\c{c}\~ao dos hiperpar\^ametros de forma autom\'atica.

% =============================================================================
% A seção de agradecimentos é opcional no arXiv
\section*{Agradecimentos}
Os autores agradecem ao suporte da CAPES e do CNPq.

% =============================================================================
\begin{thebibliography}{99}

\bibitem{arXiv2024Guide}
arXiv.org. \textit{arXiv Submission Guidelines}, 2024.
Dispon\'ivel em: \url{https://arxiv.org/help/submit}

\bibitem{arXivSubmitFAQ}
arXiv.org. \textit{Submission FAQ}, 2024.
Dispon\'ivel em: \url{https://arxiv.org/help/faq/submit}

\bibitem{Boyd2004Convex}
S. Boyd and L. Vandenberghe.
\textit{Convex Optimization}.
Cambridge University Press, 2004.

\bibitem{Goodfellow2016Deep}
I. Goodfellow, Y. Bengio, and A. Courville.
\textit{Deep Learning}.
MIT Press, 2016.

\end{thebibliography}

\end{document}
